\section*{Full Process}

\noindent\textbf{Author:} Andreas Bisiadis\\
\textbf{Affiliation:} University of West Attica

\begingroup
\sloppy

The central design constraint in this submission was that AIMO Proof Pilot was not another final-answer competition. In the preceding AIMO setting, a useful system could spend most of its budget producing many independent trajectories and then recover an integer answer by majority vote, inverse entropy weighting, or related aggregation rules. That logic does not transfer cleanly to a proof-first evaluation: a proof can be numerically correct but logically incomplete, and a large pool of mutually inconsistent attempts does not by itself identify the one solution that an olympiad grader would accept. I therefore treated the competition less as a search problem over answer strings and more as a problem of inducing proof-quality preferences into the model.

I began by selecting a base model from the restricted white-listed pool: OLMo 3 7B and 32B, OLMo 3.1 7B and 32B, Apertus 7B and 70B, and SmolLM 3B. I chose \texttt{OLMo-3.1-32B-Think}. The 7B-scale models looked too capacity-limited for proof production, especially when the evaluation permitted up to one hour per problem and therefore did not make low-latency generation the binding constraint. The 70B-scale alternatives were at the opposite extreme: they would have consumed too much of the available token-throughput budget, with no clear expected advantage over the allowed 32B class for this particular proof-writing pipeline. The 32B OLMo checkpoint was the practical middle point: it fit the intended single-accelerator and Kaggle-style inference constraints while retaining enough capacity for long-form mathematical reasoning. Within the OLMo family, the \texttt{-Think} checkpoint was especially attractive because it incorporated additional post-training beyond the base and instruction checkpoints, including later training cycles aimed at reasoning behavior. Public model-comparison tables supported this broad size and post-training tradeoff, without suggesting that the larger white-listed models would dominate strongly enough to compensate for their throughput cost.

The next issue was data. The \texttt{AI Mathematical Olympiad - Progress Prize 3} Math Corpus Prize track produced several high-quality corpora, including Crystal-Math-Preview and Astral-Math-v1, among other tool-integrated reasoning datasets. I did not use these as the main source because their construction was naturally coupled to the \texttt{AI Mathematical Olympiad - Progress Prize 3} mechanics, where the decisive output was an integer in the bounded [0, 99999] range. For Proof Pilot, I wanted data whose distribution was closer to olympiad proof writing, where the explanatory object is the main target. The MathNet release from MIT was a better match: it contained 27,817 problems sourced from national olympiads, IMO-related contests, and IMO-shortlist material. After filtering to proof-oriented, reference-bearing, text-only rows, the local builder retained 16,949 eligible proof rows, from which I created a deterministic 16-problem evaluation split and a 4,096-problem training split. This gave a compact evaluation set that was hard enough to expose proof gaps without being so difficult that a teacher model would mostly be grading noise.

My initial training hypothesis was that straightforward supervised fine-tuning would have low marginal value. The chosen checkpoint had already undergone substantial post-training, and the competition objective was not merely to imitate polished solutions but to improve the model's internal discrimination between valid and invalid reasoning. I therefore focused on Reinforcement Learning from AI Feedback (RLAIF). The intended loop was: sample a group of candidate solutions from \texttt{OLMo-3.1-32B-Think} for each MathNet problem, grade each candidate with a stronger teacher model using the Proof Pilot 0/1/6/7 mark scheme, convert that grade and auxiliary validity signals into scalar rewards, and update a LoRA adapter with a GRPO-style objective. The default group size in the implementation was 16 attempts per problem, with 32 attempts considered as a higher-budget variant.

The teacher model I planned around was \texttt{GPT-OSS-120B}. It was large enough to be a plausible proof-quality critic, compatible with the general compute plan, and already aligned with the style of recent AIMO work. The reward was not just the teacher grade: the implementation also included a context reward for nonempty, nontruncated solutions and a page-count reward so that outputs stayed near the desired proof length. In other words, the reward design tried to encode three different constraints: mathematical acceptability, completion integrity, and submission-format plausibility.

This design led to a fairly broad engineering implementation. I built deterministic MathNet preprocessing, schemas for proof records, rollout samples, reward breakdowns, and GRPO groups, a durable group queue, an OLMo ChatML tool-rollout engine, a judge prompt interface, page-counting utilities, offline and online training entrypoints, and a multi-role topology in which rollout, judge, and trainer services could run as separate nodes. The training configuration supported LoRA rank 64, alpha 128, a default learning rate of \(5 \times 10^{-6}\), BF16 training, gradient checkpointing, a 65,536-token context budget, and a queued rollout mode. The repository also contains tests for the dataset builder, reward parsing, durable queues, training dry runs, online service topology, fake inference, and the Harmony-style judge/tool loop.

The negative result was that this implementation did not mature into a successful fine-tuning run before the deadline. The failure mode was not conceptual so much as operational: the RLAIF plan required several moving pieces to work simultaneously, and debugging the pipeline through the compute containers provided by NII introduced substantial development overhead within a competition window of no more than one month. I halted the training job during the final week before the deadline so I could focus on a stable inference pipeline. The training work was therefore useful as an executable specification of the intended method, but it did not produce the model checkpoint used for the final submission. Additional time and more direct compute access would have been needed to bring the training job to submission quality.

An additional obstacle was tool-integrated reasoning. Early code paths treated Python execution through markdown code blocks as a default out-of-the-box option. However, the white-listed contestant models were not reliable enough at issuing and consuming tool calls for this to be a dependable competition-time strategy. In retrospect, tool-integrated reasoning would probably have needed to be trained before the RLAIF stage through supervised fine-tuning, rather than treated as a prompting-only capability. This shifted the final inference design away from tool use and toward a simpler proof-generation prompt.

I then tested inference-time strategies. A many-sample approach with a smaller judge selecting the best proof was rejected because the selector would itself have to solve a hard proof-quality problem, and because voting over full proofs has no natural analogue of integer majority voting. I next tried sequential refinement: the first pass generated a solution, and the second pass read the problem plus the first solution, identified gaps, and rewrote it. This was attractive because it mirrored a proof-audit workflow. In practice, on small MathNet subsets, the second pass mostly changed presentation and local wording rather than repairing the substantive mathematical gaps. It also doubled the cost for a benefit that was not visible under manual inspection. I therefore removed sequential refinement from the final path.

The final submitted system was consequently a deliberately conservative zero-shot inference pipeline using the base \texttt{OLMo-3.1-32B-Think} checkpoint. Based on observations from the recent \texttt{AI Mathematical Olympiad - Progress Prize 3} competition, where a persona-style system prompt such as IMO contestant was beneficial for the tested models, I used the same pattern in the final inference notebook. The prompt asked for an English olympiad-style proof, stable notation, explicit justification of every nontrivial step, and no private reasoning in the final answer. The local validation notebook was run on the 16-problem MathNet split with one solve pass per problem. The final local run generated 233,468 output tokens from 7,570 input tokens, for 241,038 total tokens and 11,921 seconds of elapsed generation time. This evaluation was not a formal score estimate, but it established that the single-pass pipeline could run end-to-end at the intended scale and produce correctly shaped proof submissions within the runtime envelope.

The final Kaggle-facing notebook adapted that local pipeline to the official competition paths. It used the mirrored OLMo checkpoint, launched a local vLLM OpenAI-compatible server, read the official \texttt{test.csv}, generated one proof per row, stripped hidden thinking delimiters when present, truncated each answer to the configured character budget, and wrote \texttt{submission.csv}. This was less ambitious than the intended RLAIF-trained system, but it was the stable endpoint of the process: the work clarified the reward-modeling requirements of proof evaluation, produced reusable infrastructure for a future GRPO/RLAIF run, and submitted the most reliable artifact available under the deadline.

\paragraph{Human intervention.} Human intervention was required at the main design and triage points. I chose the base model after comparing the white-listed model sizes, post-training histories, and throughput implications. I selected MathNet over \texttt{AI Mathematical Olympiad - Progress Prize 3} corpus-track datasets after judging which source better matched proof-quality evaluation. I also performed the RLAIF source-code debugging manually: failures in rollout generation, teacher scoring, queue persistence, container execution, and training launch had to be inspected from logs and reconciled against the intended GRPO data contract, rather than diagnosed by a single automated evaluation metric. That debugging shaped the final decision to treat the RLAIF implementation as an executable research scaffold rather than a submission-ready training run. I manually inspected generated MathNet proofs to determine that sequential refinement was mostly stylistic rather than corrective, and I used that assessment to remove the second pass. I also decided to abandon the partially implemented RLAIF training route for the final submission once it was clear that a reliable adapter could not be produced and validated in time. No human expert graded the full MathNet split, and no human edited the final Kaggle proofs problem by problem.

\endgroup
